\chapter{Peluang: Landasan dan Hukum Bilangan
Besar}\label{ch:b3:probability}

Tahun 2 membangun peluang pada ruang terbilang; kini teori \omterm{def:b3:measure:measure}{ukuran}
melenyapkan setiap pembatasannya. Sebuah \omterm{def:b3:probability:space}{ruang peluang} adalah
\omterm{def:b3:measure:measure}{ruang ukuran} bermassa total $1$, \omterm{def:b3:probability:space}{peubah acak} adalah
pemetaan \omterm{def:b3:lebesgue:measurable}{terukur}, \omterm{def:b3:probability:space}{nilai harapan} adalah integral Lebesgue --- dan
seketika seluruh gudang senjata analitisnya
(\cref{ch:b3:measure,ch:b3:lebesgue,ch:b3:product}) berlaku bagi
kebetulan. Bab ini memasang kamusnya, membangun
barisan tak berhingga \omterm{def:b3:probability:space}{peubah acak} yang \omterm{def:b3:probability:independence}{bebas} (pada
$\intcc01$, dari digit binernya: sebab keacakan bersembunyi di dalam
\omterm{def:b3:measure:lebesgueouter}{ukuran Lebesgue}), membuktikan lema Borel--Cantelli dan
\omterm{thm:b3:probability:zeroone}{hukum nol--satu} Kolmogorov, membereskan jenis-jenis
konvergensinya, lalu membuktikan hukum bilangan besar --- yakni
teorema yang membuat frekuensi konvergen ke peluang dan
membuat statistika mungkin. Sedangkan soal akhir pekannya menyajikan bukti
Etemadi bagi hukum kuatnya dalam bentuk $L^1$-nya yang definitif.

\section{Kamusnya}

\begin{definition}\label{def:b3:probability:space}
Sebuah \emph{ruang peluang}\index{ruang peluang} adalah
\omterm{def:b3:measure:measure}{ruang ukuran} $(\Omega, \mathcal A, \P)$ dengan $\P(\Omega) =
1$; unsur $\mathcal A$ disebut \emph{kejadian}, dan sebuah sifat
berlaku \emph{hampir pasti} (h.p.) jika kejadiannya berpeluang
$1$. Sebuah \emph{peubah acak}\index{peubah
acak} adalah pemetaan \omterm{def:b3:lebesgue:measurable}{terukur} $X \colon \Omega \to \R$ (atau ke
$\R^d$: yakni vektor acak); sedangkan \emph{distribusi}\index{distribusi sebuah
peubah acak} miliknya adalah \omterm{def:b3:measure:measure}{ukuran} peluang dorongan-maju $\P_X =
X_*\P$ pada $\R$ (\cref{exo:b3:product:9}), yang ditentukan oleh
\emph{fungsi distribusi} $F_X(t) = \P(X \leq t)$
(\cref{exo:b3:measure:3}). Lalu $X$ mempunyai \emph{\omterm{ex:b3:lebesgue:gamma}{kepadatan}} $f$ jika
$\P_X = f\,\dd\lambda$; dan ia \emph{diskret} jika $\P_X$ merupakan
kombinasi terbilang massa Dirac. Adapun
\emph{nilai harapan}\index{nilai harapan} adalah
\[
\E[X] = \int_\Omega X\,\dd\P
\qquad (X \geq 0 \text{ atau } X \in L^1(\P)),
\]
dan \emph{teorema transfer} (\cref{exo:b3:product:9})
menghitungnya di dalam distribusinya: $\E[g(X)] = \int_\R g\,\dd\P_X$ ---
yakni $= \sum g(x_k)p_k$ pada kasus diskretnya, $= \int
g(x)f(x)\dd x$ pada kasus \omterm{ex:b3:lebesgue:gamma}{kepadatannya}: yakni rumus Tahun 2, kini menjadi
teorema satu teori. Sedangkan \emph{variansnya} adalah $\V(X) =
\E[(X - \E X)^2] = \E[X^2] - (\E X)^2$ bagi $X \in L^2$.
\end{definition}

\begin{example}\label{ex:b3:probability:laws}
\omterm{def:b3:probability:space}{Distribusi} bakunya dan transformasi yang patut dicatat: Bernoulli
$\mathcal B(p)$, binomial $\mathcal B(n, p)$, geometrik,
Poisson $\mathcal P(\lambda)$ (yang diskret: tabel Tahun 2
tetap sahih); seragam pada $\intcc01$ (yakni \omterm{def:b3:measure:lebesgueouter}{ukuran Lebesgue}
itu sendiri); eksponensial $\mathcal E(\lambda)$ (berkepadatan
$\lambda\eu^{-\lambda x}\mathbf 1_{x>0}$); dan
\emph{Gauss}\index{distribusi Gauss} $\mathcal N(m, \sigma^2)$
yang berkepadatan $\frac1{\sigma\sqrt{2\pi}}\exp\bigl(-\frac{(x -
m)^2}{2\sigma^2}\bigr)$ --- yakni \omterm{ex:b3:lebesgue:gamma}{kepadatan} peluang menurut
\cref{pb:b3:lebesgue:1}, dengan rata-rata $m$ dan varians
$\sigma^2$ (yakni momen Gauss, \cref{exo:b3:product:10}).
\end{example}

\begin{proposition}[Markov dan Chebyshev]
\label{prop:b3:probability:markov}
Untuk $X \geq 0$ dan $a > 0$: $\P(X \geq a) \leq \frac{\E
X}{a}$; sedangkan untuk $X \in L^2$: $\P\bigl(\abs{X - \E X} \geq
a\bigr) \leq \frac{\V(X)}{a^2}$.\index{ketaksamaan
Markov}\index{ketaksamaan Chebyshev}
\end{proposition}

\begin{proof}
\cref{exo:b3:lebesgue:5}(a); sedangkan Chebyshev adalah Markov yang diterapkan pada
$(X - \E X)^2$.
\end{proof}

\section{Kebebasan}

\begin{definition}\label{def:b3:probability:independence}
Sub-$\sigma$-aljabar $\mathcal A_1, \dots, \mathcal A_n
\subseteq \mathcal A$ disebut \emph{bebas}\index{kebebasan}
jika $\P(A_1\cap\dots\cap A_n) = \prod\P(A_i)$ bagi setiap $A_i \in
\mathcal A_i$; kejadian disebut bebas jika $\sigma$-aljabar
$\{\varnothing, A_i, A_i^c, \Omega\}$ bebas; sedangkan \omterm{def:b3:probability:space}{peubah acak}
$X_1, \dots, X_n$ bebas jika $\sigma$-aljabar $\sigma(X_i) =
X_i^{-1}(\mathcal B(\R))$ bebas. Sebuah keluarga tak berhingga
disebut bebas jika setiap subkeluarga berhingganya bebas.
\end{definition}

\begin{theorem}\label{thm:b3:probability:independence}
Variabel $X_1, \dots, X_n$ \omterm{def:b3:probability:independence}{bebas} jika dan hanya jika \omterm{def:b3:probability:space}{distribusi} vektor
$(X_1, \dots, X_n)$ adalah \omterm{thm:b3:product:existence}{ukuran hasil kali}
$\P_{X_1}\otimes\cdots\otimes\P_{X_n}$. Dalam hal itu, bagi
$g_i \geq 0$ (atau yang hasil kalinya \omterm{def:b3:lebesgue:l1}{terintegralkan}):
\[
\E\Bigl[\prod_ig_i(X_i)\Bigr] = \prod_i\E[g_i(X_i)],
\]
khususnya $\E[XY] = \E X\,\E Y$ dan $\V(X_1 + \dots +
X_n) = \sum\V(X_i)$ bagi variabel $L^2$ yang \omterm{def:b3:probability:independence}{bebas}.
\end{theorem}

\begin{proof}
Jika $X_i$ \omterm{def:b3:probability:independence}{bebas}, maka kedua \omterm{def:b3:measure:measure}{ukuran} peluangnya
$\P_{(X_1,\dots,X_n)}$ dan $\bigotimes\P_{X_i}$ sama pada semua
hasil kali $B_1\times\dots\times B_n$ himpunan Borel --- yakni
$\pi$-sistem yang membangkitkan $\mathcal B(\R^n)$
(\cref{prop:b3:product:sections}(b)) --- sehingga sama di mana-mana
(\cref{thm:b3:measure:uniqueness}). Sebaliknya, \omterm{def:b3:probability:space}{distribusi} hasil kali
memfaktorkan setiap kejadian $\bigcap_iX_i^{-1}(B_i)$: yakni \omterm{def:b3:probability:independence}{kebebasannya}. Lalu rumus \omterm{def:b3:probability:space}{nilai harapannya} adalah
Tonelli/Fubini (\cref{thm:b3:product:tonelli}) lewat
teorema transfernya; sedangkan $\E[XY] = \E X\E Y$ adalah kasus $g_i =
\mathrm{id}$, dan menguraikan kuadratnya memberikan keaditifan
variansnya (sebab suku silangnya $\E[(X_i - \E X_i)(X_j - \E X_j)]
= 0$).
\end{proof}

\begin{theorem}[Keberadaan barisan bebas]
\label{thm:b3:probability:existence}
Pada $\bigl(\intcc01, \mathcal L, \lambda\bigr)$ ada sebuah
barisan $(U_n)_{n\geq1}$ berisi \omterm{def:b3:probability:space}{peubah acak} yang \omterm{def:b3:probability:independence}{bebas},
masing-masing seragam pada $\intcc01$. Akibatnya, bagi sembarang
\omterm{def:b3:probability:space}{distribusi} $(\mu_n)$ pada $\R$ yang ditentukan, ada $(X_n)$ \omterm{def:b3:probability:independence}{bebas} dengan
$\P_{X_n} = \mu_n$.\index{barisan bebas, keberadaan}
\end{theorem}

\begin{proof}
\emph{Digitnya.} Untuk $\omega \in \intcc01$, misalkan $(b_k(\omega))$
digit binernya (dengan $\omega = \sum b_k2^{-k}$; pilihlah
uraian yang tak berakhir dengan semua $1$ --- sebab ketaktentuannya hanya menyangkut
himpunan terbilang, jadi nol). Setiap $b_k$ merupakan \omterm{def:b3:probability:space}{peubah
acak} (sebab $\{b_k = 1\}$ gabungan berhingga selang
diadik) dan vektor $(b_1, \dots, b_m)$ mengambil setiap
nilai di $\{0,1\}^m$ pada selang diadik berpanjang $2^{-m}$:
sehingga $b_k$ merupakan Bernoulli$(\frac12)$ yang \omterm{def:b3:probability:independence}{bebas}.

\emph{Pengelompokan ulangnya.} Belahlah $\N^*$ menjadi tak berhingga banyak
himpunan tak berhingga yang saling lepas $(I_n)$ (misalnya lewat pangkat prima, atau diagonalnya);
lalu biarkanlah $(k^n_j)_j$ mencacah $I_n$ dan tetapkan
\[
U_n = \sum_{j\geq1} b_{k^n_j}\,2^{-j} .
\]
Setiap $U_n$ seragam: sebab digit binernya berupa bit adil yang
\omterm{def:b3:probability:independence}{bebas}, sehingga $\P(U_n \in [l2^{-m}, (l+1)2^{-m})) = 2^{-m}$ bagi
setiap selang diadik, sedangkan selang diadik menentukan \omterm{def:b3:probability:space}{distribusinya}
(\cref{thm:b3:measure:uniqueness}). Lalu $U_n$ saling \omterm{def:b3:probability:independence}{bebas}:
sebab semuanya fungsi blok saling lepas dari keluarga
$(b_k)$ yang \omterm{def:b3:probability:independence}{bebas} --- secara formal, kejadian $\{U_n \in D_n\}$ bagi
$D_n$ yang diadik bergantung pada berhingga banyak digit dari himpunan
yang saling lepas, sehingga memfaktor; lalu argumen $\pi$-sistemnya menaikkannya ke
semua himpunan Borel.

\emph{\omterm{def:b3:probability:space}{Distribusi} sembarang.} Ambillah $G_n(u) = \inf\{t : F_{\mu_n}(t)
\geq u\}$ (yakni \emph{fungsi kuantil} bagi fungsi
\omterm{def:b3:probability:space}{distribusi} $F_{\mu_n}$); lalu kesetaraan kuncinya $G_n(u) \leq t
\iff u \leq F_{\mu_n}(t)$ (menurut \omterm{def:b3:topology:continuity}{kekontinuan} kanan $F$ dan
kemonotonannya) menunjukkan bahwa $X_n = G_n(U_n)$ \omterm{def:b3:lebesgue:measurable}{terukur} dengan
$\P(X_n \leq t) = \P(U_n \leq F_{\mu_n}(t)) =
F_{\mu_n}(t)$: jadi berdistribusi $\mu_n$; sedangkan \omterm{def:b3:probability:independence}{kebebasannya} diwarisi
(sebab fungsi variabel yang \omterm{def:b3:probability:independence}{bebas},
\cref{exo:b3:probability:3}).
\end{proof}

\begin{example}[Masalah ulang tahun, secara jujur]
\label{ex:b3:probability:birthday}
Di antara $n$ orang yang ulang tahunnya \omterm{def:b3:probability:independence}{bebas} dan seragam atas
$N = 365$ hari, peluang bahwa semua ulang tahunnya berbeda adalah
\[
p_n = \prod_{k=1}^{n-1}\Bigl(1 - \frac kN\Bigr),
\]
lewat pensyaratan berulang (atau langsung: yakni yang menguntungkan
$N(N-1)\cdots(N - n + 1)$ atas total $N^n$, yakni argumen
pencacahan yang dijadikan ketat oleh rumus hasil kali
\omterm{def:b3:probability:independence}{kebebasannya}). Lalu dengan mengambil logaritma dan memakai $-\ln(1 - x) = x +
O(x^2)$:
\[
\ln p_n = -\frac{n(n-1)}{2N} +
O\Bigl(\frac{n^3}{N^2}\Bigr),
\qquad\text{sehingga}\qquad
p_n \approx \eu^{-n^2/2N} .
\]
Titik baliknya $p_n = \frac12$ terletak di $n \approx
\sqrt{2N\ln2} \approx 1.18\sqrt N$: untuk $N = 365$, $n = 23$
(dengan $p_{23} = 0.4927$). Ada dua pelajaran. Pertama, tumbukan di antara $n$
benda dalam $N$ kotak muncul pada skala $n \sim \sqrt N$, bukan
$n \sim N$ --- yakni \emph{penskalaan ulang tahun} yang mengatur tumbukan
hash dan ongkos $\sqrt N$ bagi serangan ulang tahun dalam
kriptografi. Kedua, perhitungannya adalah sebuah cetakan: sebab
$\binom n2$ kejadian tumbukan berpasangannya tidak \omterm{def:b3:probability:independence}{bebas}, namun
jawabannya berperilaku seolah-olah \omterm{def:b3:probability:independence}{bebas}
(sebab $\eu^{-\binom n2/N}$ persis merupakan heuristik pasangan-bebasnya)
--- yakni contoh pertama hampiran Poisson yang
dijadikan ketat pada soal akhir pekan \cref{ch:b3:clt}
(yakni ketaksamaan Le Cam).
\end{example}

\section{Borel--Cantelli dan hukum nol--satu}

\begin{theorem}[Borel--Cantelli]
\label{thm:b3:probability:borelcantelli}
Misalkan $(A_n)$ kejadian dan $\limsup A_n = \bigcap_N
\bigcup_{n\geq N}A_n$ (``$A_n$ terjadi tak berhingga
kali'').\index{lema Borel--Cantelli}
\begin{enumerate}
  \item Jika $\sum\P(A_n) < \infty$, maka $\P(\limsup A_n) =
        0$.
  \item Jika $\sum\P(A_n) = \infty$ \emph{dan $A_n$ saling
        \omterm{def:b3:probability:independence}{bebas}}, maka $\P(\limsup A_n) = 1$.
\end{enumerate}
\end{theorem}

\begin{proof}
(1) adalah \cref{exo:b3:measure:4}. (2): untuk $N \leq M$,
\omterm{def:b3:probability:independence}{kebebasan} komplemennya (\cref{exo:b3:probability:3})
memberikan
\[
\P\Bigl(\bigcap_{n=N}^{M}A_n^c\Bigr) = \prod_{n=N}^M\bigl(1
- \P(A_n)\bigr) \leq
\exp\Bigl(-\sum_{n=N}^M\P(A_n)\Bigr) \xrightarrow[M \to
\infty]{} 0
\]
(sebab $1 - x \leq \eu^{-x}$; dan deretnya divergen). Jadi
$\P\bigl(\bigcup_{n\geq N}A_n\bigr) = 1$ bagi setiap $N$, sedangkan
irisan menurun atas $N$ tetap berpeluang
$1$ (menurut \omterm{def:b3:topology:continuity}{kekontinuan} dari atas, \cref{prop:b3:measure:basics}).
\end{proof}

\begin{theorem}[Hukum nol--satu Kolmogorov]
\label{thm:b3:probability:zeroone}
Misalkan $(X_n)$ \omterm{def:b3:probability:independence}{bebas} dan $\mathcal T =
\bigcap_N\sigma(X_N, X_{N+1}, \dots)$ merupakan \emph{$\sigma$-aljabar
ekor}\index{$\sigma$-aljabar ekor} (yakni kejadian yang
tak peka terhadap berhingga banyak $X_n$ mana pun: konvergensi
$\sum X_n$, konvergensi $\frac{S_n}n$, nilai $\limsup$, \dots).
Maka setiap $T \in \mathcal T$ mempunyai $\P(T) \in \{0,
1\}$.\index{hukum nol--satu}
\end{theorem}

\begin{proof}
Tetapkanlah $N$. Kedua $\sigma$-aljabar $\sigma(X_1, \dots, X_N)$ dan
$\sigma(X_{N+1}, \dots)$ saling \omterm{def:b3:probability:independence}{bebas}: sebab kejadian yang bergantung pada
blok saling lepas memfaktor pada $\pi$-sistem pembangkitnya
(yakni silinder $\bigcap_{i\leq N}\{X_i \in B_i\}$, dan berturut-turut syarat
berhingga atas variabel berikutnya), lalu Dynkin
(\cref{thm:b3:measure:dynkin}, yang diterapkan dua kali, satu sisi setiap
kali) memperluas pemfaktorannya. Sedangkan kejadian ekor $T$ terletak di
$\sigma(X_{N+1}, \dots)$ bagi setiap $N$: sehingga $T$ \omterm{def:b3:probability:independence}{bebas} dari
setiap $\sigma(X_1, \dots, X_N)$, jadi \omterm{def:b3:probability:independence}{bebas} dari $\sigma$-aljabar
yang dibangkitkannya, $\sigma(X_1, X_2, \dots)$ (yakni Dynkin sekali lagi:
sebab gabungan $\sigma(X_1,\dots,X_N)$ merupakan $\pi$-sistem
yang membangkitkannya). Padahal $T \in \sigma(X_1, X_2, \dots)$ juga: jadi $T$
\omterm{def:b3:probability:independence}{bebas} \emph{dari dirinya sendiri}, sehingga $\P(T) = \P(T\cap T) =
\P(T)^2$: jadi $\P(T) \in \{0, 1\}$.
\end{proof}

\section{Jenis konvergensi}

\begin{definition}\label{def:b3:probability:modes}
Berlaku $X_n \to X$ \emph{hampir pasti} jika $\P(X_n \to X) = 1$;
\emph{dalam peluang} jika $\P(\abs{X_n - X} \geq \varepsilon)
\to 0$ bagi setiap $\varepsilon > 0$; dan \emph{dalam $L^p$} jika
$\E\abs{X_n - X}^p \to 0$.\index{konvergensi peubah
acak}
\end{definition}

\begin{proposition}\label{prop:b3:probability:modes}
(a) konvergensi h.p.\ mengakibatkan konvergensi dalam peluang;
(b) konvergensi $L^p$ mengakibatkan konvergensi dalam peluang;
(c) konvergensi dalam peluang mengakibatkan konvergensi h.p.\
\emph{sepanjang sebuah subbarisan};
(d) tak ada implikasi lain yang berlaku secara umum.
\end{proposition}

\begin{proof}
(a) $\P(\abs{X_n - X} \geq \varepsilon) \leq
\P\bigl(\sup_{m\geq n}\abs{X_m - X} \geq \varepsilon\bigr)
\downarrow \P\bigl(\limsup\{\abs{X_m - X} \geq
\varepsilon\}\bigr) = 0$ di bawah konvergensi h.p.\ (menurut \omterm{def:b3:topology:continuity}{kekontinuan}
dari atas; sebab kejadian limsupnya mengecualikan konvergensinya). (b)
Markov: $\P(\abs{X_n - X} \geq \varepsilon) \leq
\varepsilon^{-p}\,\E\abs{X_n - X}^p$. (c) Pilihlah $n_k$ dengan
$\P(\abs{X_{n_k} - X} \geq 2^{-k}) \leq 2^{-k}$;
lalu Borel--Cantelli (1) membuat $\abs{X_{n_k} - X} < 2^{-k}$
pada akhirnya, secara h.p. (d) Mesin tiknya (\cref{exo:b3:lp:3}) pada
$(\intcc01, \lambda)$ konvergen dalam $L^1$ dan dalam peluang
tetapi tak di mana pun secara titik demi titik; sedangkan $n\mathbf 1_{\intoo0{1/n}} \to 0$
secara h.p.\ tetapi tidak dalam $L^1$; rinciannya dan sisa
contoh tandingnya ada di \cref{exo:b3:probability:6}.
\end{proof}

\section{Hukum bilangan besar}

Sepanjang bagian ini, $(X_n)$ \omterm{def:b3:probability:independence}{bebas} dan berdistribusi sama
(\emph{i.i.d.}), dengan $S_n = X_1 + \dots + X_n$.

\begin{theorem}[Hukum lemah bilangan besar]
\label{thm:b3:probability:wlln}
Jika $X_1 \in L^2$, dengan $m = \E X_1$:
\[
\P\Bigl(\Bigl|\frac{S_n}{n} - m\Bigr| \geq
\varepsilon\Bigr) \leq
\frac{\V(X_1)}{n\,\varepsilon^2}
\xrightarrow[n\to\infty]{} 0 :
\]
sehingga $\frac{S_n}n \to m$ dalam peluang (dan dalam $L^2$).\index{hukum
bilangan besar (lemah)}
\end{theorem}

\begin{proof}
Berlaku $\E\frac{S_n}n = m$ dan $\V\bigl(\frac{S_n}n\bigr) =
\frac{n\V(X_1)}{n^2}$
(\cref{thm:b3:probability:independence}); lalu Chebyshev.
\end{proof}

\begin{theorem}[Hukum kuat bilangan besar]
\label{thm:b3:probability:slln}
Jika $X_1 \in L^1$, maka\index{hukum bilangan besar (kuat)}
\[
\frac{S_n}{n} \xrightarrow[n\to\infty]{\text{h.p.}} \E[X_1].
\]
Kita buktikan di sini di bawah hipotesis yang lebih kuat $X_1 \in
L^4$; sedangkan kasus umumnya ($L^1$: yakni bukti Etemadi) adalah
soal akhir pekannya.
\end{theorem}

\begin{proof}[Bukti di bawah $\E X_1^4 < \infty$]
Dengan memusatkannya ($X_i \mapsto X_i - m$), anggaplah $m = 0$. Lalu uraikan:
\[
\E[S_n^4] = \sum_{i,j,k,l}\E[X_iX_jX_kX_l]
= n\,\E[X_1^4] + 3n(n-1)\,\bigl(\E[X_1^2]\bigr)^2 \leq
C\,n^2 ,
\]
sebab \omterm{def:b3:probability:independence}{kebebasan} dan pemusatannya membunuh setiap suku yang memuat
faktor terpencil ($\E[X_iX_jX_kX_l] =
\E[X_i]\E[\cdots] = 0$ kecuali indeksnya berpasangan: sehingga yang
bertahan hanyalah $n$ suku $i=j=k=l$ dan $3n(n-1)$
suku dengan dua pasangan berbeda). Lalu Markov:
\[
\P\Bigl(\Bigl|\frac{S_n}n\Bigr| \geq \varepsilon\Bigr)
= \P\bigl(S_n^4 \geq n^4\varepsilon^4\bigr)
\leq \frac{Cn^2}{n^4\varepsilon^4} =
\frac{C}{\varepsilon^4n^2},
\]
yang terjumlahkan: sehingga Borel--Cantelli (1) memberikan, bagi setiap
$\varepsilon$ rasional, bahwa $\abs{S_n/n} < \varepsilon$ pada akhirnya,
secara h.p.; lalu mengirisnya atas $\varepsilon \in \Q_+^*$ (yakni terbilang banyak
kejadian berpeluang-$1$): jadi $S_n/n \to 0$ secara h.p.
\end{proof}

\begin{example}[Apa yang dibeli hukum kuatnya]
\label{ex:b3:probability:sllnapps}
(a) \emph{Frekuensinya}: bagi lemparan koin i.i.d., frekuensi
sisi gambar yang teramati konvergen secara h.p.\ ke $p$ --- yakni pembenaran
empiris bagi peluang itu sendiri.
(b) \emph{Monte Carlo}\index{metode Monte Carlo}: bagi $g \in
L^1(\intcc01)$ dan $(U_n)$ seragam i.i.d.
(\cref{thm:b3:probability:existence}),
$\frac1n\sum_{k\leq n}g(U_k) \to \int_0^1g$ secara h.p.: yakni integral
lewat pencuplikan, pada dimensi mana pun, dengan laju yang tak bergantung
dimensi $\sim n^{-1/2}$ yang dipertajam pada \cref{ch:b3:clt}.
(c) \emph{Bilangan normal}: hampir setiap bilangan real mempunyai, pada
uraian binernya, frekuensi asimtotik $\frac12$ bagi angka satu
(terapkanlah hukum kuatnya pada variabel digit
\cref{thm:b3:probability:existence}) --- yakni teorema Borel, sebuah
pernyataan tentang bilangan \emph{sehari}-hari yang dibuktikan lewat \omterm{def:b3:measure:measure}{ukuran}:
\cref{pb:b3:probability:1} melengkapinya pada semua basis.
\end{example}

\begin{method}\label{met:b3:probability:toolkit}
Urutan kerja bagi pernyataan asimtotik tentang barisan
acak: (1) \emph{Apakah kejadiannya kejadian ekor?} Jika ya, peluangnya
$0$ atau $1$
(\cref{thm:b3:probability:zeroone}) dan tinggal
diputuskan yang mana. (2) \emph{Untuk membuktikan pernyataan h.p.}:
Borel--Cantelli --- yakni peluang yang terjumlahkan bagi
kejadian yang ``buruk'', lewat batas bertipe Markov/Chebyshev atas
momen apa pun yang ada; sedangkan \omterm{def:b3:probability:independence}{kebebasannya} hanya diperlukan bagi
arah konversnya. (3) \emph{Subbarisan + apitan}: buktikanlah
konvergensinya sepanjang subbarisan yang terkelola, lalu kendalikanlah
ayunan di antaranya lewat kemonotonan atau ketaksamaan
maksimal --- yakni kerangka bukti Etemadi. (4) Untuk
limit \omterm{def:b3:probability:space}{distribusinya}, tunggulah \cref{ch:b3:clt}.
\end{method}

\section{Latihan}

\begin{exercise}[$\star$]\label{exo:b3:probability:1}
(a) Misalkan $X$ berfungsi \omterm{def:b3:probability:space}{distribusi} $F$ yang \omterm{def:b3:topology:continuity}{kontinu} dan naik
sejati. Tunjukkanlah bahwa $F(X)$ seragam pada $\intcc01$, dan
bahwa $G(U) \sim F$ bagi $U$ yang seragam, dengan $G = F^{-1}$: yakni simulasi
lewat pembalikan.
(b) Hitunglah fungsi \omterm{def:b3:probability:space}{distribusi} dan \omterm{ex:b3:lebesgue:gamma}{kepadatan} $X^2$
bagi $X$ yang seragam pada $\intcc{-1}1$, dan bagi $-\frac1\lambda\ln
U$ dengan $U$ yang seragam pada $\intoo01$.
\end{exercise}

\begin{exercise}[$\star$]\label{exo:b3:probability:2}
(a) Hitunglah rata-rata dan varians \omterm{def:b3:probability:space}{distribusi} Poisson $\mathcal
P(\lambda)$ dan geometrik lewat teorema transfernya.
(b) Tunjukkanlah bahwa \omterm{def:b3:probability:space}{peubah acak} positif $T$ dengan $\P(T > t)
> 0$ bagi setiap $t$ memenuhi sifat \emph{tanpa ingatan} $\P(T > t + s \mid
T > t) = \P(T > s)$ bagi setiap $s, t \geq 0$ jika dan hanya jika $T$
eksponensial. \emph{(Sebab fungsi kesintasannya memenuhi persamaan
fungsional Cauchy; sedangkan kemonotonannya menggantikan \omterm{def:b3:topology:continuity}{kekontinuannya}.)}
\end{exercise}

\begin{exercise}[$\star\star$]\label{exo:b3:probability:3}
(a) Tunjukkanlah bahwa jika $X_1, \dots, X_n$ \omterm{def:b3:probability:independence}{bebas} dan $f_i$
fungsi Borel, maka $f_i(X_i)$ \omterm{def:b3:probability:independence}{bebas}.
(b) Tunjukkanlah bahwa kejadian $A_1, \dots, A_n$ \omterm{def:b3:probability:independence}{bebas} jika dan hanya jika
komplemennya \omterm{def:b3:probability:independence}{bebas}, jika dan hanya jika indikator $\mathbf 1_{A_i}$
merupakan \omterm{def:b3:probability:space}{peubah acak} yang \omterm{def:b3:probability:independence}{bebas}.
(c) (\omterm{def:b3:probability:independence}{Bebas} berpasangan itu lebih lemah) Dua koin adil: $A =$ yang pertama
gambar, $B =$ yang kedua gambar, $C =$ keduanya sama. Tunjukkanlah bahwa
$A, B, C$ \omterm{def:b3:probability:independence}{bebas} berpasangan tetapi tidak \omterm{def:b3:probability:independence}{bebas}.
\end{exercise}

\begin{exercise}[$\star\star$]\label{exo:b3:probability:4}
(a) (Monyet tak berhingga) Barisan i.i.d.\ ketukan tombol seragam
pada abjad berhingga secara h.p.\ memuat setiap teks berhingga
tak berhingga kali: buktikanlah dengan Borel--Cantelli (2) pada
blok yang saling lepas.
(b) (Rentetan) Bagi bit adil i.i.d., misalkan $R_n$ panjang
rentetan angka satu yang mulai di posisi $n$. Tunjukkanlah bahwa secara h.p.\
$R_n \geq (1+\varepsilon)\log_2n$ hanya berhingga kali, sedangkan $R_n
\geq \log_2 n$ tak berhingga kali \emph{(lewat kedua paruh
Borel--Cantelli; dan untuk yang kedua, berpindahlah ke blok saling lepas demi
memperoleh \omterm{def:b3:probability:independence}{kebebasan})}: sehingga rentetan terpanjang pada $n$ digit pertamanya
tumbuh seperti $\log_2n$.
\end{exercise}

\begin{exercise}[$\star\star$]\label{exo:b3:probability:5}
Misalkan $(X_n)$ \omterm{def:b3:probability:independence}{bebas}.
(a) Tunjukkanlah bahwa jari-jari konvergensi $\sum X_n z^n$ merupakan
konstanta h.p.\ (mungkin $0$ atau $\infty$).
(b) Tunjukkanlah bahwa $\P(\sum X_n \text{ konvergen}) \in \{0, 1\}$
dan $\P(S_n/n \to m) \in \{0,1\}$.
(c) Berikanlah kejadian tentang $(X_n)$ yang \emph{bukan} kejadian
ekor, lalu periksalah bahwa \omterm{thm:b3:probability:zeroone}{hukum nol--satunya} bisa gagal padanya.
\end{exercise}

\begin{exercise}[$\star\star$]\label{exo:b3:probability:6}
Pada $(\intcc01, \lambda)$, tunjukkanlah --- beserta buktinya --- \omterm{def:b3:probability:space}{peubah
acak} sedemikian sehingga:
(a) $X_n \to 0$ dalam peluang dan dalam setiap $L^p$, tetapi
tak di mana pun secara h.p.;
(b) $X_n \to 0$ secara h.p.\ tetapi tak dalam $L^p$ mana pun;
(c) $X_n \to 0$ dalam $L^1$ tetapi tak dalam $L^2$;
(d) lalu tunjukkanlah: jika $X_n \to X$ dalam peluang dan $\abs{X_n}
\leq Y \in L^1$, maka $X_n \to X$ dalam $L^1$
\emph{(lewat subbarisan + konvergensi terdominasi +
kiat subsubbarisannya)}.
\end{exercise}

\begin{exercise}[$\star\star$]\label{exo:b3:probability:7}
Sebuah jajak pendapat menaksir proporsi $p$ yang tak diketahui lewat
frekuensi empiris $\hat p_n$ atas $n$ penarikan yang \omterm{def:b3:probability:independence}{bebas}.
(a) Chebyshev: tunjukkanlah $\P(\abs{\hat p_n - p} \geq \varepsilon)
\leq \frac1{4n\varepsilon^2}$ (pakailah $p(1-p) \leq \frac14$).
(b) Berapa banyak penarikan yang menjamin galat $\leq 3\%$ dengan
peluang $\geq 95\%$ menurut batas ini? (Jawaban sebenarnya,
lewat \cref{ch:b3:clt}, kira-kira $1070$: jadi Chebyshev jujur
tetapi kasar.)
\end{exercise}

\begin{exercise}[$\star\star\star$]\label{exo:b3:probability:8}
(Bernstein) Untuk $f \in \mathcal C(\intcc01)$ definisikanlah
polinomial Bernstein $B_nf(x) =
\sum_{k=0}^n\binom nkx^k(1-x)^{n-k}f\bigl(\frac kn\bigr)$.
(a) Kenalilah $B_nf(x) = \E\bigl[f\bigl(\frac
{S_n}n\bigr)\bigr]$ bagi $S_n$ yang binomial $\mathcal B(n, x)$.
(b) Buktikanlah $B_nf \to f$ \emph{secara seragam} pada $\intcc01$: belahlah
pada $\{\abs{\frac{S_n}n - x} \leq \delta\}$ dan
komplemennya, dengan memakai \omterm{def:b3:topology:continuity}{kekontinuan} seragam dan Chebyshev beserta
batas seragam $\V(\frac{S_n}n) \leq \frac1{4n}$.
(c) Simpulkanlah: yakni bukti kedua yang bersifat peluang bagi
teorema hampiran Weierstrass
(\cref{cor:b3:complete:weierstrass}), dengan laju eksplisit
$\norm{B_nf - f}_\infty \leq \frac32\,\omega_f(n^{-1/2})$
bagi modulus \omterm{def:b3:topology:continuity}{kekontinuan} $\omega_f$ --- buktikanlah setidaknya
bentuk $O(\omega_f(n^{-1/2}))$-nya.
\end{exercise}

\begin{exercise}[$\star\star\star$]\label{exo:b3:probability:9}
(Pengumpul kupon) Kartu dari $n$ jenis ditarik secara seragam
dengan pengembalian; misalkan $T_n$ banyaknya penarikan sampai semua
jenisnya terlihat.
(a) Tulislah $T_n = \sum_{k=1}^{n}\tau_k$ dengan $\tau_k$ yang
geometrik berparameter $\frac{n - k + 1}n$, dan $\tau_k$ yang
\omterm{def:b3:probability:independence}{bebas}, lalu simpulkanlah $\E T_n = n\,H_n \sim n\ln n$
(dengan $H_n$ bilangan harmoniknya) dan $\V(T_n) \leq
\frac{\pi^2}6n^2$.
(b) Chebyshev: $\frac{T_n}{n\ln n} \to 1$ dalam peluang.
(c) Pertajamlah dengan Borel--Cantelli: tunjukkanlah secara langsung $\P(T_n >
\beta n\ln n) \leq n^{1 - \beta}$ bagi $\beta > 1$
\emph{(lewat batas gabungan atas kejadian bahwa suatu jenis terlewat
setelah $\beta n\ln n$ penarikan, dengan memakai $1 - x \leq \eu^{-x}$)},
lalu simpulkanlah bahwa sepanjang $n = 2^m$, secara h.p.\ $T_n \leq \beta n\ln
n$ pada akhirnya, bagi setiap $\beta > 2$.
\end{exercise}

\begin{exercise}[$\star\star$]\label{exo:b3:probability:10}
Dengan memakai konstruksi digitnya
(\cref{thm:b3:probability:existence}):
(a) periksalah lewat perhitungan langsung bahwa $U = \sum b_{2k}2^{-k}$
(yakni digit berindeks genap milik $\omega$ yang seragam) bersifat seragam dan
\omterm{def:b3:probability:independence}{bebas} dari $V = \sum b_{2k-1}2^{-k}$;
(b) simpulkanlah adanya bijeksi \omterm{def:b3:lebesgue:measurable}{terukur} sampai himpunan nol antara
$\intcc01$ dan $\intcc01^2$ yang mengawetkan \omterm{def:b3:measure:measure}{ukuran}, lalu berilah komentar:
bahwa satu bilangan acak seragam memuat dua (bahkan terbilang banyak)
bilangan \omterm{def:b3:probability:independence}{bebas} --- bandingkanlah dengan kurva Peano
(\cref{pb:b3:topology:1}), yang mencapai kesurjektifan tetapi
tidak pengawetan \omterm{def:b3:measure:measure}{ukuran} ataupun keinjektifan.
\end{exercise}

\begin{exercise}[$\star\star$]\label{exo:b3:probability:11}
(Rekor) Misalkan $(X_n)_{n\geq1}$ i.i.d.\ dengan fungsi
\omterm{def:b3:probability:space}{distribusi} yang \omterm{def:b3:topology:continuity}{kontinu}, dan katakanlah sebuah \emph{rekor} terjadi pada
waktu $n$ jika $X_n > \max(X_1, \dots, X_{n-1})$ (dan waktu $1$
adalah rekor). Misalkan $R_n$ indikator rekornya.
(a) Tunjukkanlah $\P(R_n = 1) = \frac1n$ \emph{(sebab menurut kesetangkupannya, masing-masing dari
$n!$ pengurutan $X_1, \dots, X_n$ sama mungkinnya
sedangkan seri berpeluang $0$)}.
(b) Tunjukkanlah bahwa $R_n$ bersifat \emph{\omterm{def:b3:probability:independence}{bebas}} \emph{(cacahlah
pengurutan yang cocok dengan posisi rekor yang ditentukan, atau
berargumenlah bahwa urutan relatif $X_1, \dots, X_{n-1}$ bersifat
\omterm{def:b3:probability:independence}{bebas} dari peringkat $X_n$ di antaranya)}.
(c) Simpulkanlah dari Borel--Cantelli
(\cref{thm:b3:probability:borelcantelli}, kedua paruhnya) bahwa
tak berhingga banyak rekor terjadi secara h.p., tetapi rekor pada
waktu berurutan $n, n+1$ terjadi tak berhingga kali dengan
peluang --- putuskanlah yang mana! --- lalu hitunglah
$\sum_n\P(R_n = 1, R_{n+1} = 1)$.
\end{exercise}

\begin{exercise}[$\star\star$]\label{exo:b3:probability:12}
(Rentetan gambar terpanjang) Lemparkanlah koin adil tak berhingga kali, lalu
misalkan $L_n$ panjang rentetan terpanjang sisi gambar yang
berurutan di dalam $n$ lemparan pertamanya.
(a) Tunjukkanlah bahwa bagi setiap $\varepsilon > 0$, secara h.p.\ $L_n \leq
(1 + \varepsilon)\log_2n$ pada akhirnya \emph{(sebab peluang
bahwa suatu rentetan berpanjang $\ell$ mulai di antara $n$ lemparan
pertamanya paling banyak $n2^{-\ell}$; lalu Borel--Cantelli sepanjang $n =
2^k$)}.
(b) Tunjukkanlah bahwa secara h.p.\ $L_n \geq (1 - \varepsilon)\log_2n$
pada akhirnya \emph{(potonglah $n$ lemparan pertamanya menjadi
$\lfloor n/\ell\rfloor$ blok saling lepas berpanjang $\ell =
\lceil(1 - \varepsilon)\log_2n\rceil$; blok itu saling
\omterm{def:b3:probability:independence}{bebas}, masing-masing seluruhnya gambar dengan peluang $2^{-\ell}$,
sedangkan peluang bahwa tak satu pun seluruhnya gambar paling banyak
$\exp(-n2^{-\ell}/\ell)$; lalu jumlahkanlah sepanjang $n = 2^k$ lagi)}.
(c) Simpulkanlah $\frac{L_n}{\log_2n} \to 1$ secara h.p.: bahwa dalam sejuta
lemparan adil kita patut menantikan rentetan sekitar $20$ gambar ---
dan sekumpulan data tanpa itu boleh jadi dikarang.
\end{exercise}

\section{Soal: bukti Etemadi bagi hukum kuatnya}

\begin{problem}[{Soal akhir pekan --- hukum kuat bilangan
besar bagi variabel i.i.d.\ yang terintegralkan}]
\label{pb:b3:probability:1}
Hukum kuat Kolmogorov --- $\frac{S_n}n \to \E X_1$ secara h.p.\
bagi $X_n \in L^1$ yang i.i.d. --- lama hanya mempunyai bukti yang berbelit;
lalu pada 1981 N.~Etemadi menemukan satu bukti yang sangat hemat, tanpa memakai
apa pun di luar bab ini (bahkan dengan melemahkan \omterm{def:b3:probability:independence}{kebebasannya}
menjadi \omterm{def:b3:probability:independence}{kebebasan} berpasangan). Kita ikuti buktinya. Misalkan $(X_n)$
\omterm{def:b3:probability:independence}{bebas} berpasangan, berdistribusi identik, dan \omterm{def:b3:lebesgue:l1}{terintegralkan};
dengan $m = \E X_1$, $S_n = X_1 + \dots + X_n$.

\medskip
\noindent\textbf{Bagian I --- Reduksinya.}
\begin{enumerate}
  \item Tunjukkanlah bahwa cukup ditangani $X_n \geq 0$
        \emph{(belahlah $X_n = X_n^+ - X_n^-$: lalu periksalah bahwa kedua
        paruhnya kembali \omterm{def:b3:probability:independence}{bebas} berpasangan, i.i.d., dan
        \omterm{def:b3:lebesgue:l1}{terintegralkan})}. Selanjutnya anggaplah $X_n \geq 0$.
  \item (Pemenggalan) Misalkan $Y_n = X_n\,\mathbf 1_{X_n \leq n}$
        dan $S_n^* = Y_1 + \dots + Y_n$. Tunjukkanlah
        \[
        \sum_{n\geq1}\P(X_n \neq Y_n) =
        \sum_{n\geq1}\P(X_1 > n) \leq \E[X_1] < \infty
        \]
        (\cref{exo:b3:product:3}), lalu simpulkanlah lewat
        Borel--Cantelli bahwa $\frac{S_n - S_n^*}{n} \to 0$
        secara h.p.: sehingga cukup dibuktikan $\frac{S^*_n}n \to m$
        secara h.p.
  \item Tunjukkanlah $\E Y_n = \E\bigl[X_1\mathbf 1_{X_1\leq
        n}\bigr] \to m$ (menurut konvergensi monoton), sehingga
        $\frac1n\sum_{k\leq n}\E Y_k \to m$ (Cesàro): jadi
        cukup dibuktikan $\frac{S_n^* - \E S_n^*}{n} \to 0$
        secara h.p.
\end{enumerate}

\medskip
\noindent\textbf{Bagian II --- Taksiran variansnya.}
\begin{enumerate}[resume]
  \item Tunjukkanlah
        \[
        \V(Y_n) \leq \E[Y_n^2] = \E\bigl[X_1^2\,\mathbf
        1_{X_1 \leq n}\bigr]
        \]
        lalu, dengan memakai kue berlapisnya
        (\cref{prop:b3:product:layercake}), batas kuncinya
        \[
        \sum_{n\geq1}\frac{\V(Y_n)}{n^2}
        \leq \sum_{n\geq1}\frac1{n^2}\,
        \E\bigl[X_1^2\mathbf 1_{X_1\leq n}\bigr]
        \leq C\,\E[X_1] < \infty
        \]
        \emph{(tukarlah jumlahnya dengan \omterm{def:b3:probability:space}{nilai harapannya} ---
        yakni Tonelli bagi deret --- lalu batasilah $\sum_{n \geq
        x}\frac1{n^2} \leq \frac2{\max(x,1)}$ bagi taksiran
        dalamnya $x^2\sum_{n\geq x}n^{-2} \leq 2x$)}.
\end{enumerate}

\medskip
\noindent\textbf{Bagian III --- Konvergensi sepanjang subbarisan
geometrik.} Tetapkanlah $\alpha > 1$ lalu misalkan $k_j =
\lfloor\alpha^j\rfloor$.
\begin{enumerate}[resume]
  \item Dengan memakai \omterm{def:b3:probability:independence}{kebebasan} berpasangannya (sebab variansnya menjumlah,
        \cref{thm:b3:probability:independence} --- periksalah bahwa
        keaditifan variansnya hanya memerlukan \omterm{def:b3:probability:independence}{kebebasan}
        berpasangan) dan Chebyshev, tunjukkanlah bagi setiap
        $\varepsilon > 0$:
        \[
        \sum_{j\geq1}\P\Bigl(\Bigl|
        \frac{S^*_{k_j} - \E S^*_{k_j}}{k_j}\Bigr| \geq
        \varepsilon\Bigr)
        \leq
        \frac1{\varepsilon^2}\sum_{j\geq1}\frac1{k_j^2}
        \sum_{n\leq k_j}\V(Y_n)
        = \frac1{\varepsilon^2}\sum_{n\geq1}\V(Y_n)
        \sum_{j\,:\,k_j\geq n}\frac1{k_j^2} .
        \]
  \item Tunjukkanlah $\sum_{j : k_j \geq n}k_j^{-2} \leq
        \frac{C_\alpha}{n^2}$ \emph{(lewat deret geometrik;
        dan waspadalah pada fungsi lantainya: yakni kehati-hatian bertipe $k_j \geq \frac{\alpha^j}2$ bagi
        $\alpha^j \geq 2$)}, lalu simpulkanlah dengan
        pertanyaan 4 dan Borel--Cantelli:
        \[
        \frac{S^*_{k_j} - \E S^*_{k_j}}{k_j}
        \xrightarrow[j\to\infty]{\text{h.p.}} 0,
        \qquad\text{sehingga}\qquad
        \frac{S^*_{k_j}}{k_j} \to m \ \text{h.p.}
        \]
\end{enumerate}

\medskip
\noindent\textbf{Bagian IV --- Apitan dan kesimpulannya.}
\begin{enumerate}[resume]
  \item Untuk $k_j \leq n \leq k_{j+1}$, pakailah kemonotonan
        $S^*_n$ (sebab sukunya tak negatif!) untuk menunjukkan
        \[
        \frac{k_j}{k_{j+1}}\,\frac{S^*_{k_j}}{k_j}
        \;\leq\; \frac{S^*_n}{n} \;\leq\;
        \frac{k_{j+1}}{k_j}\,\frac{S^*_{k_{j+1}}}{k_{j+1}},
        \]
        lalu simpulkanlah, secara h.p.:
        \[
        \frac m\alpha \leq \liminf\frac{S^*_n}n \leq
        \limsup\frac{S^*_n}n \leq \alpha\,m .
        \]
  \item Biarkanlah $\alpha \downarrow 1$ sepanjang sebuah barisan lalu
        simpulkanlah $\frac{S_n^*}n \to m$ secara h.p., sehingga (menurut Bagian I)
        \emph{hukum kuat bilangan besar}:
        \[
        \boxed{\ \frac{S_n}{n}
        \xrightarrow[n\to\infty]{\text{h.p.}} \E[X_1].\ }
        \]
  \item Di manakah persisnya \omterm{def:b3:probability:independence}{kebebasan} berpasangan (alih-alih
        \omterm{def:b3:probability:independence}{kebebasan} penuh) sudah mencukupi? Daftarkanlah ketiga tempat
        yang hipotesis bertipe \omterm{def:b3:probability:independence}{kebebasan} dipanggil di sana.
\end{enumerate}

\medskip
\noindent\textbf{Bagian V --- Dividennya.}
\begin{enumerate}[resume]
  \item (Bilangan normal Borel) Tunjukkanlah bahwa menurut
        $\lambda$, hampir setiap $x \in \intcc01$ bersifat
        \emph{normal pada setiap basis} $b \geq 2$: yakni setiap digit
        $0, \dots, b-1$ muncul dengan frekuensi asimtotik
        $\frac1b$ \emph{(tetapkanlah $b$ dan sebuah digit, lalu terapkan
        hukum kuatnya pada variabel indikatornya --- benarkanlah
        bahwa digit basis-$b$ sebuah variabel seragam bersifat
        seragam i.i.d.\ pada $\{0,\dots,b-1\}$ seperti pada
        \cref{thm:b3:probability:existence} --- lalu
        iriskanlah terbilang banyak kejadian
        berpeluang satu itu)}. Tunjukkanlah satu bilangan tak normal yang eksplisit,
        lalu renungkanlah: bahwa teoremanya menegaskan kenormalan hampir
        semua bilangan, padahal membuktikan kenormalan $\sqrt2$ atau
        $\pi$ tetap terbuka.
  \item (\omterm{ex:b3:probability:sllnapps}{Monte Carlo}, terjamin) Benarkanlah selengkapnya
        metode \cref{ex:b3:probability:sllnapps}(b) bagi
        $g \in L^1(\intcc01^d)$: bangunlah cuplikan seragam
        i.i.d.\ pada $\intcc01^d$ dari
        \cref{thm:b3:probability:existence} dan
        \cref{exo:b3:probability:10}, lalu nyatakanlah apa yang
        diberikan hukum kuatnya.
\end{enumerate}

\medskip
\noindent\textbf{Bagian VI --- Apa yang dibeli \omterm{def:b3:probability:independence}{kebebasan} penuh:
ketaksamaan maksimal dan deret acak.} Etemadi hanya
membelanjakan \omterm{def:b3:probability:independence}{kebebasan} berpasangan; sedangkan bagian sisanya memanfaatkan
versi penuhnya (yang bersama). Misalkan $(Z_n)$ variabel $L^2$
terpusat yang \omterm{def:b3:probability:independence}{bebas} dan $S_k = Z_1 + \dots + Z_k$ (yakni notasi
baru, yang tak berkaitan dengan $X_n$ di atas).
\begin{enumerate}[resume]
  \item (Ketaksamaan maksimal Kolmogorov) Untuk $\varepsilon
        > 0$ buktikanlah
        \[
        \P\Bigl(\max_{1\leq k\leq n}\abs{S_k} \geq
        \varepsilon\Bigr) \;\leq\;
        \frac1{\varepsilon^2}\sum_{k=1}^n\V(Z_k) :
        \]
        Ongkos Chebyshev membeli maksimumnya \emph{(partisikanlah
        kejadiannya menurut indeks pertama $k$ yang
        $\abs{S_k} \geq \varepsilon$; lalu pada keping itu tulislah
        $S_n^2 \geq S_k^2 + 2S_k(S_n - S_k)$ dan pakailah
        \omterm{def:b3:probability:independence}{kebebasan} koalisi $(Z_1, \dots, Z_k)$
        dan $(Z_{k+1}, \dots, Z_n)$,
        \cref{thm:b3:probability:independence})}. Tunjukkanlah
        langkah yang \omterm{def:b3:probability:independence}{kebebasan} berpasangannya tak lagi
        mencukupi di sana.
  \item (Teorema satu deret Khinchin--Kolmogorov) Simpulkanlah:
        bahwa jika $\sum_n\V(Z_n) < \infty$, maka $\sum_nZ_n$
        konvergen hampir pasti \emph{(tunjukkanlah bahwa secara h.p.\
        jumlah parsialnya membentuk barisan Cauchy: biarkanlah $m \to
        \infty$ pada ketaksamaan maksimal yang diterapkan pada
        $Z_{N+1}, \dots, Z_{N+m}$, lalu biarkan $N \to
        \infty$)}.
  \item (Deret Rademacher) Misalkan $(\varepsilon_n)$
        tanda i.i.d., dengan $\P(\varepsilon_n = \pm1) =
        \frac12$ (\cref{thm:b3:probability:existence}), dan
        misalkan $(x_n)$ bilangan real. Tunjukkanlah bahwa
        $\sum_nx_n\varepsilon_n$ konvergen secara h.p.\ begitu
        $\sum_nx_n^2 < \infty$; tunjukkanlah pula bahwa, apa pun
        $(x_n)$-nya, peluang bahwa
        $\sum_nx_n\varepsilon_n$ konvergen adalah $0$ atau $1$
        (\cref{thm:b3:probability:zeroone}).
  \item Konversnya, secara elementer. Tetapkanlah $T_n = \sum_{k\leq
        n}x_k\varepsilon_k$ dan $s_n^2 = \sum_{k\leq
        n}x_k^2$, lalu andaikanlah $s_n \to \infty$.
        (a) Buktikanlah \emph{ketaksamaan Paley--Zygmund}: bahwa bagi
        $Z \geq 0$ dengan $\E Z^2 < \infty$ dan $0 < \theta <
        1$,
        \[
        \P\bigl(Z > \theta\,\E Z\bigr) \;\geq\; (1 -
        \theta)^2\,\frac{(\E Z)^2}{\E Z^2}
        \]
        \emph{(belahlah $\E Z$ pada aras $\theta\E Z$ lalu
        terapkanlah Cauchy--Schwarz pada keping atasnya)}.
        (b) Tunjukkanlah $\E T_n^4 \leq 3s_n^4$.
        (c) Simpulkanlah $\P\bigl(\abs{T_n} > \frac{s_n}2\bigr)
        \geq \frac3{16}$ lalu simpulkanlah bahwa
        $\sum_nx_n\varepsilon_n$ divergen secara h.p.; sehingga
        dikotominya
        \[
        \sum_nx_n\varepsilon_n\ \text{konvergen h.p.}
        \iff \sum_nx_n^2 < \infty .
        \]
  \item (Deret harmonik acak) Simpulkanlah bahwa
        $\sum_n\frac{\varepsilon_n}{n^s}$ konvergen secara h.p.\
        jika dan hanya jika $s > \frac12$. Untuk $\frac12 < s \leq
        1$ deretnya konvergen secara h.p.\ padahal
        $\sum_nn^{-s} = \infty$: sebab tanda acak menghasilkan
        peniadaan sekuat akar kuadrat --- bandingkanlah dengan
        deret berselang-seling $\sum_n\frac{(-1)^n}{n^s}$,
        yang konvergen bagi \emph{setiap} $s > 0$.
\end{enumerate}

\medskip
\noindent\textbf{Bagian VII --- Pemusatan: ketaksamaan
Hoeffding.} Hukum kuatnya mengatakan $\frac{S_n}n \to m$;
sedangkan ketaksamaan pemusatan mengatakan betapa tak mungkinnya sebuah simpangan
\emph{pada setiap $n$ yang tetap}.
\begin{enumerate}[resume]
  \item (Lema Hoeffding)
        (a) Tunjukkanlah $\cosh\lambda \leq \eu^{\lambda^2/2}$ bagi
        setiap $\lambda \in \R$, dengan membandingkan kedua deretnya
        suku demi suku.
        (b) Misalkan $Z$ terpusat dengan $a \leq Z \leq b$,
        $a < b$. Tunjukkanlah
        \[
        \E\,\eu^{\lambda Z} \leq
        \exp\Bigl(\frac{\lambda^2(b - a)^2}8\Bigr)
        \]
        \emph{(batasilah $\eu^{\lambda z}$ pada $\intcc ab$ oleh
        talinya, ambillah \omterm{def:b3:probability:space}{nilai harapannya}, lalu pelajarilah $\varphi(t)
        = -pt + \log(1 - p + p\eu^t)$ dengan $p =
        \frac{-a}{b-a}$ dan $t = \lambda(b - a)$: tunjukkanlah
        $\varphi(0) = \varphi'(0) = 0$ dan $\varphi''
        \leq \frac14$)}.
  \item (Ketaksamaan Hoeffding) Misalkan $X_1, \dots, X_n$
        \omterm{def:b3:probability:independence}{bebas} dengan $a_i \leq X_i \leq b_i$ dan $S_n =
        X_1 + \dots + X_n$. Buktikanlah, bagi $t > 0$,
        \[
        \P\bigl(S_n - \E S_n \geq t\bigr) \leq
        \exp\Bigl(\frac{-2t^2}{\sum_{i=1}^n(b_i -
        a_i)^2}\Bigr),
        \]
        beserta batas yang sama bagi ekor bawahnya
        \emph{(lewat Chebyshev eksponensial: batasilah
        $\E\,\eu^{\lambda(S_n - \E S_n)}$ dengan memakai
        \omterm{def:b3:probability:independence}{kebebasannya} dan pertanyaan 17, lalu optimumkanlah atas
        $\lambda > 0$)}.
  \item (Hukum kuatnya, kasus terbatas, dengan laju) Misalkan
        $X_i$ i.i.d.\ dengan nilai di $\intcc ab$ dan $m
        = \E X_1$. Tunjukkanlah
        \[
        \P\Bigl(\Bigl|\frac{S_n}n - m\Bigr| \geq
        \varepsilon\Bigr) \leq
        2\exp\Bigl(\frac{-2n\varepsilon^2}{(b - a)^2}\Bigr)
        \]
        lalu perolehlah kembali $\frac{S_n}n \to m$ secara h.p.\ lewat
        Borel--Cantelli: yakni bukti kedua bagi hukum kuatnya
        bagi variabel terbatas --- tanpa pemenggalan, dengan laju
        eksponensial pada setiap $n$ berhingga, tetapi dengan suku terbatas
        dan \omterm{def:b3:probability:independence}{kebebasan} penuh. Bandingkanlah
        hipotesisnya dengan hipotesis Etemadi.
  \item (\omterm{ex:b3:probability:sllnapps}{Monte Carlo}, terjamin pada $n$ yang tetap) Misalkan $g
        \colon \intcc01^d \to \intcc01$ \omterm{def:b3:lebesgue:measurable}{terukur} dan
        $(U_k)$ cuplikan seragam i.i.d.\ pertanyaan 11.
        Diberikan $\varepsilon, \delta > 0$, tunjukkanlah
        \[
        n \geq \frac{\log(2/\delta)}{2\varepsilon^2}
        \implies
        \P\Bigl(\Bigl|\frac1n\sum_{k=1}^ng(U_k) -
        \int g\,\dd\lambda_d\Bigr| \geq \varepsilon\Bigr)
        \leq \delta,
        \]
        lalu hitunglah ambangnya bagi $\varepsilon =
        \delta = 10^{-2}$. Batasnya tak melibatkan $d$:
        bandingkanlah dengan pertanyaan 11 dan dengan kisi
        deterministiknya.
\end{enumerate}

\medskip
\noindent\textbf{Bagian VIII --- Seberapa besar jalan acak?
Menuju logaritma teriterasi.} Misalkan $S_n = \varepsilon_1 +
\dots + \varepsilon_n$ jalan acak sederhana yang dibangun dari
tanda adil i.i.d.
\begin{enumerate}[resume]
  \item (Ekor sub-Gauss) Tunjukkanlah $\E\,\eu^{\lambda S_n} =
        (\cosh\lambda)^n \leq \eu^{n\lambda^2/2}$ lalu
        simpulkanlah, bagi $x > 0$,
        \[
        \P(S_n \geq x) \leq \eu^{-x^2/(2n)},
        \qquad
        \P(\abs{S_n} \geq x) \leq 2\,\eu^{-x^2/(2n)} .
        \]
  \item Simpulkanlah, lewat Borel--Cantelli,
        \[
        \limsup_{n\to\infty}\frac{\abs{S_n}}
        {\sqrt{2n\log n}} \leq 1 \quad\text{h.p.}
        \]
        \emph{(bagi $\eta > 0$, jumlahkanlah batas ekornya di $x =
        (1 + \eta)\sqrt{2n\log n}$, lalu iriskanlah atas
        $\eta = \frac1p$)}. Khususnya jalan acaknya hidup pada
        skala teorema limit pusatnya $\sqrt n$ sampai faktor logaritmik
        --- jauh di bawah batas kasarnya $\abs{S_n} \leq n$.
  \item Sepanjang subbarisan penggandaan $n_j = 2^j$, tunjukkanlah
        \[
        \limsup_{j\to\infty}\frac{S_{n_j}}
        {\sqrt{2n_j\log\log n_j}} \leq 1 \quad\text{h.p.},
        \]
        lalu renungkanlah: bahwa \emph{hukum logaritma
        teriterasi} (Khinchin; dan Hartman--Wintner bagi suku
        $L^2$ terpusat yang umum) menyatakan bahwa
        \[
        \limsup_{n\to\infty}\frac{S_n}
        {\sqrt{2n\log\log n}} = 1 \quad\text{h.p.}
        \]
        Jelaskanlah dengan tepat apa yang memisahkan taksiran
        subbarisan yang baru dibuktikan itu dari paruh atas
        pernyataan ini (sebab orang harus mengendalikan $\max_{n_j \leq n \leq
        n_{j+1}}S_n$ di dalam setiap bloknya, dan itu menuntut
        ketaksamaan maksimal pada skala \emph{eksponensial})
        lalu periksalah secara kuantitatif bahwa ketaksamaan
        pertanyaan 12 terlalu lemah untuk maksud itu. Sedangkan paruh
        bawahnya bersandar pada lema Borel--Cantelli kedua
        yang diterapkan pada blok yang \omterm{def:b3:probability:independence}{bebas}; kedua paruhnya
        merupakan bahan Tahun 3 yang jujur bagi kuliah peluang
        tersendiri.

  \item (Simpangan seragam atas kelas berhingga) Misalkan $A_1,
        \dots, A_N$ kejadian pada sebuah percobaan yang dapat diulang,
        lalu taksirlah setiap peluangnya lewat frekuensi
        empirisnya $\hat p_i$ atas $n$ pengulangan i.i.d.
        Dengan menggabungkan ketaksamaan Hoeffding dan batas gabungan,
        tunjukkanlah
        \[
        \P\Bigl(\max_{i\leq N}\,\abs{\hat p_i - \P(A_i)} >
        \varepsilon\Bigr) \;\leq\; 2N\,\eu^{-2n\varepsilon^2},
        \]
        lalu simpulkanlah kaidah \omterm{def:b3:measure:measure}{ukuran} cuplikannya: bahwa $n \geq
        \frac{\ln(2N/\delta)}{2\varepsilon^2}$ menjamin
        seluruh $N$ taksirannya sekaligus
        akurat-$\varepsilon$ dengan peluang $\geq 1 -
        \delta$. Hitunglah $n$ bagi $N = 10^6$, $\varepsilon =
        0.01$, $\delta = 0.05$: yakni harga logaritmik bagi
        keseragamannya.
  \item (Jendela harmonik acaknya) Dengan menggabungkan kedua
        paruh teori deret acaknya, tunjukkanlah bahwa bagi
        tanda i.i.d.\ $(\varepsilon_n)$ deret
        $\sum_n\frac{\varepsilon_n}{n^\alpha}$ konvergen
        secara h.p.\ jika $\alpha > \frac12$ dan divergen secara h.p.\ jika
        $\alpha \leq \frac12$; lalu pertentangkanlah dengan konvergensi
        mutlaknya (yang menuntut $\alpha > 1$): sebab pada
        jendela $\alpha \in \intoc{\frac12}1$, konvergensinya
        merupakan gejala yang sungguh bersifat peluang ---
        yakni peniadaan, bukan \omterm{def:b3:measure:measure}{ukuran}.

\end{enumerate}
\end{problem}
